\chapter{Introduction}
\label{ch:introduction}
Research into small-scale biological phenomena, e.g., observing sub-cellular processes, understanding genetics, and tracking molecular machines, has for many decades driven advances in microscopy. One of these was the use of fluorescence as a contrast mechanism, i.e. the attachment of fluorescent agents to biological structures, exciting them with laser-produced light of a specific wavelength, and observing the emitted light which has a different wavelength, to form an image of the underlying structure.

To capture the 3D complexity of biological processes without losing spatial resolution, fluorescence microscopy was transformed into a 3D imaging mechanism via the use of sequential depth scanning, at the cost of both temporal resolution and increased risk of photo-toxicity due to extended exposure times.

Depth capture and its costs are not the only challenges to overcome in fluorescence microscopy setups, since light that travels through a complex sample and a sequence of different optical components accumulates distortions to its wavefront, which we collectively call optical aberrations. Software-based solutions for aberration correction were developed, usually requiring point spread function (PSF) estimation and deconvolution in post-processing, which is computationally expensive, time-consuming and prone to artifacts and noise amplification. Hardware-based solutions were also proposed, which involve the use of Shack--Hartmann wavefront sensors to estimate the aberrations and a correction process using adaptive optics. These methods have been very successful, but they require expensive hardware like deformable mirror (DM) devices to operate in real-time, and can usually only apply spatially invariant corrections --- a narrow capability compared to the complexity of aberrations in biological samples.

A different capture strategy addresses the depth problem directly while simultaneously encoding the aberrations compactly: light-field microscopy (LFM) manages to capture spatial information together with the angular information by placing a microlens array in the native image plane (NIP), allowing a single-capture 3D reconstruction without the need for scanning while aberrations are encoded as relative tilts in different parts of the aperture. Though innovative and promising, LFM entailed two main challenges: a spatially uneven sampling pattern, and a high computational cost of performing estimation and deconvolution. To address these, researchers created the Fourier light-field microscope (FLFM), which moves the microlens array from the NIP to the pupil-conjugate plane, thus creating a better aliased sampling pattern and a cheaper deconvolution process \cite{guo2019fourierLightField}. While the new design offers advantages, a central tradeoff is carried over from LFM: larger and fewer microlenses allow better spatial resolution but reduce angular resolution, and vice versa. 

Placing a microlens array in a pupil-conjugate plane is geometrically the arrangement of the Shack--Hartmann wavefront sensor~\cite{platt2001shackHartmannHistory} used in adaptive optics: each lenslet samples a sub-aperture of the pupil, and a local wavefront tilt over that sub-aperture displaces the image it forms. In an FLFM, however, each lenslet images the extended scene rather than a guide star. A single capture therefore contains both the aberrations, encoded as relative shifts between views, and the images needed to correct them.

In this thesis, we present a new approach to the problem of aberration correction in light-field microscopy. Our system is based on an FLFM architecture that is aimed at capturing an aberrated 2D scene, leveraging the angular information of the light-field to estimate the aberrations, which appear as simple shifts in first-order approximation, and correcting them using a global-then-local digital registration. The estimation-correction process is done entirely in post-processing without the need for active components, scanning, or deconvolution operations, making it simple and fast. It also allows for a relatively simple and inexpensive hardware setup, while still being able to image and correct scenes with spatially varying aberrations with a high-NA objective and extended depth of field, in a single capture.
